\section{Data, Preparation, and Descriptive Analysis}
\label{sec:data_prep}

% Arc: everything about the data prep, front and center — then preliminary
% descriptive analysis (content TBD), ending with per-bar OLS on HAR as the
% shared benchmark every later section is measured against.

\subsection{Data}
\label{sec:dp_data}

The raw data are 30-minute bars for the S\&P~500 complex, gridded to a
regular 30-minute frequency and spanning January 2005 through April 2024.
The panel is built on a near-24-hour clock covering pre-market, regular, and
after-hours sessions, running Sunday 18:30 through Friday 20:00, so each
trading day contributes $48$ bars --- roughly $34$ overnight or
extended-session bars and roughly $14$ regular-session bars. This is an
easily missed feature of the design: the majority of bars are \emph{not}
regular-session bars, and any statement about ``intraday'' structure in this
panel is a statement about a 24-hour cycle, not about the
$09{:}30$--$16{:}00$ window alone. Periods in which no venue is open are
removed by a fixed calendar filter: all bars on Friday after $20{:}00$, all
of Saturday, and all of Sunday before $18{:}30$. No other calendar
exclusions are applied; holidays and half-days remain in the panel and are
handled by the availability machinery of Section~\ref{sec:dp_prep} rather
than by deletion. The full regular grid after this filter contains
$246{,}057$ rows, from 2005-01-02 to 2024-04-30.

The forecasting target is realized \emph{variance}, measured as the sum of
squared 5-minute log-returns within each 30-minute bar:
\begin{equation}
    RV_{d,j} = \sum_{i=1}^{M} r_{d,j,i}^{2}, \qquad M = 6,
    \label{eq:dp_rv}
\end{equation}
where $r_{d,j,i}$ is the $i$-th 5-minute log-return inside the $j$-th
30-minute bar of day $d$. We index bars sequentially as $RV_t$ and forecast
one bar ahead, $RV_{t+1}$. We use ``realized variance'' throughout for the
target itself and reserve ``volatility'' for its square root, on which the
models are estimated (Section~\ref{sec:dp_prep}). This is the standard
realized-measure construction, sampled at the intraday frequency for which
the heterogeneous autoregressive model of \citet{Corsi2009} was designed but
pushed one rung finer: the HAR literature's canonical unit is the trading
day, whereas here both the target and the lag ladder live at 30-minute
resolution.

\paragraph{The evaluation panel.}
Feature construction consumes a burn-in of $3{,}125$ bars, the length of
the longest HAR lag, which reduces the usable panel to $242{,}934$ bars.
The walk-forward protocol then consumes a further $24{,}000$ bars as the
initial training window, leaving $n = 218{,}934$ evaluated bars. Every arm
reported in this paper is scored on those identical rows, which is what
licenses the bar-for-bar paired comparisons used throughout.

\paragraph{Exogenous information.}
Beyond the target's own history we use $41$ exogenous channels, organised
into thematic buckets that partition the columns exactly
($7+12+6+6+3+3+4 = 41$):
\begin{itemize}
    \item \textbf{Moments} ($7$): own-asset higher-order return moments ---
      bipower variation, realized semi-variance, autocovariance, absolute,
      cubic, and quartic returns.
    \item \textbf{Liquidity} ($12$): trading volume, turnover (total,
      buy-initiated, sell-initiated), effective and quoted spreads, number
      of trades.
    \item \textbf{Market, equal-weighted} ($6$) and \textbf{value-weighted}
      ($6$): cross-sectional stock-factor aggregates of return moments,
      bipower variation, semi-variance, and absolute returns.
    \item \textbf{Sentiment} ($3$): StockTwits attention, sentiment score,
      and sentiment count.
    \item \textbf{Implied volatility} ($3$): VIX, VVIX, VIX3M.
    \item \textbf{Volatility demand} ($4$): options order-flow indicators
      (SPX open/close, all options open/close, open only).
\end{itemize}
The buckets are of very unequal richness --- twelve columns of liquidity
against three of sentiment --- so per-bucket comparisons confound
information content with dimensionality unless read per feature. And
availability is \emph{structurally} unequal: the cross-sectional aggregates
are observed on $138{,}357$ bars, the options order-flow bucket on only
$54{,}062$, and the VIX family does not exist for the earliest part of the
sample at all. Handling this by row deletion would score each bucket on its
own, different, sub-sample; Section~\ref{sec:dp_prep} describes the
impute-and-indicate construction that removes the confound instead.

\pending{feature accounting: a per-bucket table giving channels, the lag
set applied to each, derived columns, and the reconciliation of the three
circulating counts (41 channels, 359-column exogenous block, 529-column
design) --- the stated expansion rule of six rolling means per channel gives
$41 \times 6 = 246$ and does not reconcile them; no per-column claim should
be read as final until this table exists}

\subsection{Preparation}
\label{sec:dp_prep}
% The full prep pipeline as first-class material: overnight fills, diurnal
% adjustment, winsorized sqrt target, rolling robust scaling, availability
% indicators (the stale-data lesson folded in as prep discipline, not as the
% paper's centerpiece).

We treat the preparation pipeline as first-class material rather than as
housekeeping, for a reason this project learned at some cost: measured on
this panel, the representation of missing and stale data is a larger effect
than any model-class, basis, or penalty decision we report. At any given
bar, $36\%$ of the exogenous design is not a measurement but an echo of an
earlier one --- on the sparsest channels the figure is $78\%$ --- and a model
handed those columns without honest bookkeeping is fitting the fill policy,
not the market. This subsection states the pipeline and, where a stage
exists because a defect taught us it must, says so.

\paragraph{Overnight fills and availability indicators.}
The $41$ exogenous channels are not all published on the near-24-hour grid.
Most are US cash- or extended-session equity quantities, each printing only
inside its own session; averaged across channels, $0.638$ of panel rows
carry a genuine print and the remaining $0.362$ must be filled. We
forward-fill each channel within its availability window, substitute an
in-distribution neutral value where it is structurally absent, and
accompany every value column with a binary availability indicator. The
indicator is computed \emph{before} the fill, from the raw feed --- the
ordering is load-bearing. An earlier build of this pipeline computed it
after the fill, so that every forward-filled bar was labelled observed; when
the volatility-demand feed stopped publishing on 2023-08-31, the unlimited
fill carried its final August print across the following eight months
($8{,}453$ bars, $3.44\%$ of the sample) with the indicator reading
$\approx 1$ throughout, and three further feeds (\texttt{vix},
\texttt{vix3m}, \texttt{vvix}) terminated the same way on 2024-02-12. We now
gate every built cache on a single check: for each channel, the mean of the
constructed availability indicator must equal the fraction of rows on which
the raw feed actually printed. The two quantities are estimates of the same
number; any gap is the bug or one like it. On the uncorrected panel the
check fails for $31$ of $41$ channels. It requires no model, no held-out
data, and no threshold tuning, and we recommend it as standard for any
mixed-frequency panel in which exogenous series are filled.

The indicator does double duty in estimation. The linear coefficient on it
absorbs any level shift between the pre- and post-availability eras, and the
value column contributes only where it exists; every bucket, including
implied volatility and volatility demand, therefore lives on the same
$218{,}934$-row index, which is precisely what makes the paired
significance tests of the later sections exact.

\paragraph{Stage 1: diurnal adjustment.}
Intraday volatility has a pronounced, largely deterministic time-of-day
profile. Let $s(t) \in \{1,\dots,48\}$ index the time-of-day slot of bar
$t$ and let $\mathcal{W}_t = \{j : j < t,\ s(j)=s(t)\}$ be the set of the
most recent $W = 20$ observations sharing that slot, with at least $5$
required before the baseline is valid. For non-negative series the baseline
is the same-slot rolling mean and the adjustment divides by it,
\begin{equation}
    \bar{RV}_{t\mid s} = \frac{1}{|\mathcal{W}_t|}\sum_{j \in \mathcal{W}_t} RV_j,
    \qquad
    \widetilde{RV}_t = \frac{RV_t}{\bar{RV}_{t\mid s}},
    \label{eq:dp_diurnal}
\end{equation}
while signed series (e.g.\ autocovariance) are divided by the same-slot
rolling standard deviation instead. The restriction $j < t$, not
$j \le t$, is what makes the adjustment non-anticipative. Warm-up bars pass
through with baseline $1$ and are consumed by the burn-in, so they never
enter the evaluation panel.

\paragraph{Stage 2: the winsorized square-root target.}
Each variable class receives a variance-stabilising map chosen from its
support and tail behaviour rather than fitted: $x \mapsto \sqrt{x}$ for
squared-return and volatility-like measures, including $RV$ itself, bipower
variation, turnover, and effective spread; $x \mapsto
\mathrm{sign}(x)\sqrt{|x|}$ for autocovariance; $x \mapsto x^{1/3}$ and
$x \mapsto x^{1/4}$ for third and fourth moments; and $x \mapsto \log x$
for the implied-volatility family. Models are estimated on the adjusted
$\sqrt{RV}$ scale. After transformation, each series is clipped to rolling
5th and 95th percentile bounds estimated on the trailing
$W_{\mathrm{clip}} = 240$ observations (about five trading days) ending at
$t-1$. Winsorization is lossy and is not inverted at evaluation time; only
the diurnal division and the square root are. Because
$\mathbb{E}[\sqrt{X}]^2 \neq \mathbb{E}[X]$, naively squaring a forecast of
$\sqrt{RV}$ is downward biased, and we apply Duan's smearing correction
\citep{Duan1983}, adding the in-sample residual variance before squaring and
then undoing the diurnal division:
\begin{equation}
    \widehat{RV}^{\mathrm{raw}}_t
    = \left(\hat{f}_t^{\,2} + \hat{\sigma}^2_{\varepsilon}\right)\cdot
      \bar{RV}^{\mathrm{baseline}}_{t\mid s},
    \qquad
    \hat{\sigma}^2_{\varepsilon} = \frac{1}{N}\sum_i (y_i - \hat{y}_i)^2.
    \label{eq:dp_duan}
\end{equation}

\paragraph{HAR features.}
The target's own history enters through trailing rolling means of the
adjusted series on a base-5 geometric ladder,
\begin{equation}
    \mathcal{L} = \{1,\ 5,\ 25,\ 125,\ 625,\ 3125\},
    \qquad
    \overline{RV}^{(k)}_t = \frac{1}{k}\sum_{i=0}^{k-1}\widetilde{RV}_{t-1-i},
    \label{eq:dp_har}
\end{equation}
spanning roughly $30$ minutes, $2.5$ hours, half a trading day, $2.6$ days,
$13$ days, and $65$ trading days; the one-bar shift keeps contemporaneous
information out of the feature. This is the multi-scale device of
\citet{Corsi2009} transposed from the daily/weekly/monthly triple to a
six-rung intraday ladder. The same ladder is applied to each included
exogenous channel, contributing $|\mathcal{L}| = 6$ rolling-mean columns per
channel.

\paragraph{Rolling robust scaling.}
Input features to the linear estimators are standardised by a rolling
robust scaler --- median and interquartile range estimated from the current
training buffer, updated at every step. Two refinements to this scaler are,
again, defects turned into discipline. First, location and scale are
estimated from \emph{observed} rows only: estimating them over
forward-filled values lets a frozen series collapse its own rolling IQR
toward zero, turning the standardising division into a division by nearly
nothing --- on the uncorrected panel the stale volatility-demand columns
reached $2{,}260$ rolling IQRs not because the series moved but because the
denominator vanished. Second, masked estimation has its own failure mode at
the opposite end of a feed's life: in a feed's opening weeks there are few
observed rows, and a quiet opening fortnight can set the scale for
everything that follows. We therefore require $512$ observations before a
scale is trusted and floor the IQR at $1\%$ of the largest scale the
channel has shown. The repair is an interaction --- honest indicators,
masked estimation, and neutral substitution each fail when shipped alone
--- and composed they take the worst channel from $2{,}260.4$ to $96.2$
rolling IQRs. On the affected October 2023 window the difference between
the uncorrected and repaired panel was QLIKE $6.37170$ against $0.12391$;
every genuine modelling effect reported in this paper lies between
$0.0005$ and $0.004$. That ratio is why preparation leads this paper's
exposition rather than trailing it.

All stages are strictly causal: every quantity dated $t$ uses only
information available strictly before $t$.

\subsection{Descriptive Analysis}
\label{sec:dp_descriptive}

\decide{what preliminary descriptive analysis goes here --- candidates for
which material already exists in the sources:
(i) the persistence result --- a pure persistence forecast beats both direct
kNN and spectral-embedding kNN on the six HAR features in the prototype
walk-forward (MSE $0.076277$ vs $0.087887$ and $0.089581$, against a
moving-average baseline at $0.107461$), establishing that the conditional
mean is overwhelmingly the recent level;
(ii) the diurnal profile and its residual --- Stage 1 removes $77\%$ of the
diurnal variance, the $23\%$ leak retains about $48\%$ of the original
amplitude but only about $4\%$ of residual noise (the profile figure itself
does not yet exist as an artifact);
(iii) the clock-anchored persistence sign-flip --- corr(\texttt{har\_ma\_5},
residual) is $+0.05$ to $+0.13$ through the continuous session but
$-0.085$ at the open and $-0.085/-0.21/-0.14/-0.03$ across hours 16--19,
the motivation for the linear-vs-tree question;
(iv) design-matrix geometry --- $p=529$ columns, participation ratio
$\approx 6.8$, eigenvalue decay $\approx i^{-3.13}$, numerical rank $317$,
PC1 carrying $\approx 0\%$ of target signal: the dense-but-weak fact;
(v) crisis-window conduct --- the 2008 and 2020 figures showing the HAR
baseline tracking the level regime but undershooting peak bars, motivating
QLIKE;
(vi) an era overview / raw-$RV$ summary-statistics table by session segment
(no such table exists yet).
Note that (i)--(v) were computed for the reference paper's panel and would
need re-verification on the current frozen panel before inclusion.}

\subsection{The OLS--HAR Benchmark}
\label{sec:dp_benchmark}
% Endpoint of the section: per-bar OLS on the six HAR features. The number the
% rest of the paper improves on.

Every result in the remainder of this paper is referenced to a single fixed
comparator: ordinary least squares on the six HAR features of
Equation~\eqref{eq:dp_har} plus calendar dummies, with no exogenous
variables, refit at every bar. Forecasts are produced by a strictly causal
rolling-origin procedure with a training window of $T_{\mathrm{train}} =
24{,}000$ bars (approximately $500$ trading days at $48$ bars per day),
sliding by one bar per prediction; the per-bar refit is computationally
feasible because the sliding-window normal equations admit exact rank-1
up- and down-dating, and the benchmark is obtained from the ridge path at
$\alpha = 0$, hence exact OLS. Forecasts are mapped back to raw variance
through Equation~\eqref{eq:dp_duan} and scored by QLIKE
\citep{Patton2011},
\begin{equation}
    L_{\mathrm{QLIKE}}
    = \frac{1}{N}\sum_{t=1}^{N}\left[\, r_t - \log r_t - 1 \,\right],
    \qquad
    r_t = \frac{RV^{\mathrm{raw}}_t}{\widehat{RV}^{\mathrm{raw}}_t},
    \label{eq:dp_qlike}
\end{equation}
on the full $218{,}934$-bar evaluation panel. QLIKE is scale-invariant in
the level of volatility, strictly consistent for the conditional variance,
and penalises under-prediction more heavily than over-prediction --- the
economically relevant asymmetry for this target. It is, however, not
comparable across preprocessing regimes: on the diurnally adjusted panel of
the reference study the same class of models scores several tenths of a
QLIKE unit lower than on a raw, unadjusted panel, so every comparison in
this paper stays strictly within one regime and every table is labelled
with the panel it was computed on.

\pending{the benchmark table: per-bar OLS-on-HAR QLIKE and secondary
metrics (out-of-sample $R^2$, Mincer--Zarnowitz $\alpha$/$\beta$, MSE, MAE)
on the current frozen panel. The reference paper's incumbent scored QLIKE
$0.13415$, $R^2_{\mathrm{OOS}}$ $0.694$, MZ $\beta$ $0.930$ on its
$218{,}934$-bar panel (\texttt{incumbent\_metrics.json}); those values must
be re-confirmed or re-run before this table is frozen, since the present
paper's panel and target conventions may differ}

Differences in QLIKE at the fourth decimal place are not self-evidently
meaningful, so every comparison against this benchmark is accompanied by a
Diebold--Mariano test on the per-bar loss differential
\citep{DieboldMariano1995} with a Newey--West long-run variance and the
small-sample correction of \citet{HarveyLeybourneNewbold1997}; where
several arms are compared simultaneously we report the Model Confidence Set
of \citet{HansenLundeNason2011}. Because every arm lives on the identical
row index, all of these are exact paired bar-for-bar comparisons.
